\chapter{Semaine 3 }

L'objectif de cette est de poursuivre l'analyse et le traitement de la 
base de données. 
Par la suite nous allons débuter l'Implémentation de modèle de classification 
afin de débuter les prédictions des modèles.


\section{Séléction des features}

Un fois les déscripteurs calculés, nous devons choisir les quels conserver pour obtenir 
les meilleurs résultats possible. 
Dans cet objectif nous avons explorer la méthode RFE(recursive feature elimination). 
Étant donné un estimateur externe qui pondère les caractéristiques, l'objectif 
de l'élimination récursive de caractéristiques (RFE) est de sélectionner les 
caractéristiques en considérant successivement des ensembles de plus en plus petits.
La variant que nous avons utilisé est RFECV, CV signifie simplement validation 
croisée (Cross validation). Le nombre de caractéristiques sélectionnées est ajusté 
automatiquement en adaptant un sélecteur RFE aux différentes plis de validation 
croisée. Dans notre cas il est important de ne pas oublier que les données 
possèdent un ordre chronologique, ainsi nous devons utiliser un \textit{train/test split} 
adéquat, c'est pour cela que nous avons utilisé \textit{TimeSeriesSplit}.  